\documentclass{article}

\usepackage[spanish]{babel}
\usepackage{amsmath,amssymb}

\title{Certamen 3 Cálculo}
\author{Benjamín Rivas Beltrán}
\date{\today}

\begin{document}

\maketitle

\section{Suma áreas cuadrados}

Dado que el primer cuadrado tiene lado 1 y cada cuadrado tiene lado igual a la mitad del lado del cuadrado anterior, el lado el $n$-ésmio cuadrado es de $l_n = ( \frac{1}{2})^{n - 1}$. Luego, el área del $n$-ésmio cuadrado es:
\begin{align*}
A_n &= (l_n)^2 \\ 
&= \left(\left(\frac{1}{2}\right)^{n - 1}\right)^2 \\
&= \left(\frac{1}{4}\right)^{n - 1} \\ 
\end{align*}

La suma de las áreas de todos los cuadrados es:
\begin{align*}
S &= \sum_{n = 1}^{\infty} A_n \\
&= \sum_{n = 1}^{\infty} \left(\frac{1}{4}\right)^{n - 1}
\end{align*}

La cual es una serie geométrica de razón $\frac{1}{4} < 1$, que se sabe que converge a:
\begin{align*}
S &= \frac{1}{1 - \frac{1}{4}} \\
&= \frac{1}{\frac{3}{4}} \\
&= \boxed{\frac{4}{3}}
\end{align*}

Por tanto, la suma de las áreas de todos los cuadrados es $\frac{4}{3}$.

\section{Serie por sumas parciales}

\[
S = \sum_{n = 1}^{\infty} \frac{1}{(n + 2)(n + 3)}
\]

Reescribiendo el término de la serie con fracciones parciales:

\begin{align*}
\frac{1}{(n + 2)(n + 3)} &= \frac{A}{n + 2} + \frac{B}{n + 3}     \ \ / * (n + 2)(n + 3) \\
1 &= A(n + 3) + B(n + 2)
\end{align*}

Reemplazando para $n = -2$:
\begin{align*}
1 &= A (3 - 2) + B(2 - 2) \\
1 &= A 
\end{align*}

Reemplazando para $n = -3$:
\begin{align*}
1 &= A (3 - 3) + B(2 -32) \\
1 &= -B \\
-1 &= B 
\end{align*}

Reemplazando en la serie:

\begin{align*}
S &= \sum_{n = 1}^{\infty} \left[\frac{A}{n + 2} + \frac{B}{n + 3}\right] \\
&= \sum_{n = 1}^{\infty} \left[\frac{1}{n + 2} - \frac{1}{n + 3}\right]
\end{align*}

Sea $a_n = \frac{1}{n + 2}$, entonces:
\[
S = \sum_{n = 1}^{\infty} [a_n - a_{n + 1}]
\]

Sea $N \in \mathbb{N}$ y:
\[
S_N = \sum_{n = 1}^{N} [a_n - a_{n + 1}]
\]

Expandiendo:
\begin{align*}
S_N &= (a_1 - a_2) + (a_2 - a_3) + \ldots + (a_N - a_{N + 1}) \\
&= a_1 - a_2 + a_2 - a_3 + \ldots + a_N - a_{N + 1}
\end{align*}

Note que todos los términos se cancelan, salvo el primero y el último, por tanto:
\begin{align*}
S_N &= a_1 - a_{N + 1} \\
S_N &= \frac{1}{1 + 2} - \frac{1}{N + 1 + 2} \\
S_N &= \frac{1}{3} - \frac{1}{N + 3}
\end{align*}

El valor de la serie es:
\begin{align*}
S &= \lim_{N \to \infty} S_N \\
&= \lim_{N \to \infty} \left[\frac{1}{3} - \frac{1}{N + 3} \right] \\
&= \lim_{N \to \infty} \frac{1}{3} - \lim_{N \to \infty} \frac{1}{N + 3} \\
&= \frac{1}{3} - 0 \\
&= \boxed{\frac{1}{3}}
\end{align*}

\[
\therefore S = \sum_{n = 1}^{\infty} \frac{1}{(n + 2)(n + 3)} = \frac{1}{3}
\]

\section{Convergencia de series}

\subsection{$\sum_{n = 1}^{\infty} \frac{(-3)^{2n + 1}}{5^n}$}

Reescribiendo la serie:

\begin{align*}
S &= \sum_{n = 1}^{\infty} \frac{(-3)^{2n + 1}}{5^n} \\
&= \sum_{n = 1}^{\infty} \frac{(-3)^{2n} \cdot (-3) }{5^n} \\
&= \sum_{n = 1}^{\infty} \frac{\left((-3)^2\right)^n \cdot (-3) }{5^n} \\
&= -3 \sum_{n = 1}^{\infty} \frac{9^n}{5^n}\\
&= -3 \sum_{n = 1}^{\infty} \left( \frac{9}{5} \right)^n
\end{align*}

La cual es una serie geométrica de razón $r = \frac{9}{5} > 1$, que se sabe que diverge.

Por lo tanto, la serie:

\[
\sum_{n = 1}^{\infty} \frac{(-3)^{2n + 1}}{5^n}
\]

Es divergente.

\subsection{$\sum_{n = 1}^{\infty} \left( \frac{n}{n + 1} \right)^{-2n} $}

Se sabe que para que la serie $\sum_{n = 1}^{\infty} a_n$ sea convergente, se debe cumplir que $\lim_{n \to \infty} a_n = 0$. 

Desarrollando el límite:
\begin{align*}
L = \lim_{n \to \infty} a_n &= \lim_{n \to \infty} \left( \frac|{n}{n + 1} \right)^{-2n} \\
&= \lim_{n \to \infty} \left( \frac{n + 1}{n} \right)^{2n} \\
&= \lim_{n \to \infty} \left( 1 + \frac{1}{n} \right)^{2n} \\
&= \left(\lim_{n \to \infty} \left( 1 + \frac{1}{n} \right)^{n} \right)^2 \\
\end{align*}

Por definición del número de Euler:
\[
\lim_{n \to \infty} \left( 1 + \frac{1}{n} \right)^{n} = e
\]

Luego:
\[
L = \lim_{n \to \infty} a_n = e^2 \neq 0
\]

Por lo tanto, la serie:

\[
\sum_{n = 1}^{\infty} \left( \frac{n}{n + 1} \right)^{-2n} 
\]

Es divergente.

\subsection{$\sum_{n = 1}^{\infty} n^2 e^{-n^3}$}

Sea $f: [3, \infty) \to \mathbb{R} $ una función dada por:
\[
f(x) = x^2 e^{-n^3}
\]

Note que la función es continua, pues es el producto de dos funciones continuas. Además, como $e^x > 0$ y $x^2 > 0$ para todo $x > 0$, la función es positiva en todo $[3, \infty)$.

Derivando:
\begin{align*}
f'(x) &= 2x e^{-x^3} + x^2 e^{-x^3} (-3x^2) \\
&= e^{-x^3} (2x - 3x^4)
\end{align*}

Note que para $x \geq 3$:
\begin{align*}
x^3 &\geq 27 \ / \cdot -3 \\
-3 x^3 &\leq -81 \\
2 - 3 x^3 &\leq 2 - 81 \\
2 - 3 x^3 &\leq -79 \ / \cdot x > 0 \\
2x - 3 x^4 &\leq -79x < 0 \ / \cdot e^{-x^3} > 0 \\
e^{-x^3} (2x - 3 x^4) &< 0 \\
f'(x) < 0
\end{align*}

Por tanto, $f$ es decreciente en todo $[3, \infty)$.

Con eso, podemos aplicar el criterio de la integral, el cual establece que la serie converge si y solo si la integral:

\[
\int_{3}^{\infty} f(x)dx
\]

Converge.

Desarrollando la integral:
\begin{align*}
\int_{3}^{\infty} f(x)dx &= \int_{3}^{\infty} x^2 e^{-x^3}dx \\
&= \lim_{t \to \infty} \int_{3}^{t} x^2 e^{-x^3}dx
\end{align*}

Cambio de variable $u = x^3$, diferenciando: $du = 3x^2 dx \leftarrow x^2 dx = \frac{du}{3}$

\begin{align*}
\lim_{t \to \infty} \int_{3}^{t^3} e^{-u} \frac{du}{3} &= \lim_{t \to \infty}-\frac{1}{3} e^{-u} \big|_0^{t^3} \\ 
&= -\frac{1}{3} \lim_{t \to \infty} (e^{-t^3} - 1) \\ 
&= -\frac{1}{3} (-1)\\ 
&= -\frac{1}{3}
\end{align*}

Por lo tanto, la serie:

\[
\sum_{n = 1}^{\infty} n^2 e^{-n^3}
\]

Es convergente.

\section{Series de potencias}

\subsection{$\sum_{n = 1}^{\infty} \frac{x^n}{n!}$}

Por el criterio de la razón, la serie $\sum_{n = 1}^{\infty} a_n$ converge si $\lim_{n \to \infty} \frac{a_{n + 1}}{a_n} < 1$

Desarrollando:
\begin{align*}
\lim_{n \to \infty} \frac{a_{n + 1}}{a_n} &= \lim_{n \to \infty} \frac{\frac{x^{n + 1}}{(n + 1)!}}{\frac{x^n}{n!}} \\
&= \lim_{n \to \infty} \frac{x}{n + 1} = 0 < 1, \forall x \in \mathbb{R} \\
\end{align*}

Por tanto, la serie converge para todo x.

\subsection{$f(x) = e^{3x}$}

Derivando:
\begin{align*}
f'(x) = 3e^{3x} \\
f''(x) = 9e^{3x} \\
f^{(n)}(x) = 3^n e^{3x} \\
f^{(n)}(0) = 3^n
\end{align*}

\[
f(x) = \sum_{n = 1}^{\infty} \frac{3^n x^n}{n!} = \sum_{n = 1}^{\infty} \frac{3^n x^n}{n!}
\]

\end{document}